% Generated draft: requires mathematical and source review.
\documentclass[11pt]{article}
\usepackage[margin=25mm]{geometry}
\usepackage{fontspec}
\setmainfont{Latin Modern Roman}
\usepackage{amsmath,amssymb,mathtools,unicode-math}
\setmathfont{Latin Modern Math}
\usepackage{enumitem,needspace,xurl,hyperref}
\setlength{\parindent}{0pt}
\setlength{\parskip}{5pt}
\setlength{\emergencystretch}{3em}
\newcommand{\N}{\mathbb{N}}
\newcommand{\Z}{\mathbb{Z}}
\newcommand{\Q}{\mathbb{Q}}
\newcommand{\R}{\mathbb{R}}
\newcommand{\C}{\mathbb{C}}
\newcommand{\F}{\mathbb{F}}
\newcommand{\A}{\mathbb{A}}
\newcommand{\PP}{\mathbb{P}}
\newcommand{\E}{\mathbb{E}}
\newcommand{\eps}{\varepsilon}
\newcommand{\dd}{\,\mathrm{d}}
\begin{document}
\section*{UnsolvedMath ID: 30001893}
\textbf{Source problem code:} catalogue entry
\subsection{Definitions and mathematical statement}
A set \(C\subseteq\R^3\) is convex if \((1-t)x+ty\in C\) for all \(x,y\in C\) and all \(t\in[0,1]\).

The supplied problem asks: for each integer \(n>3\), determine the dimension of the space of all partitions of \(\R^3\) into exactly \(n\) convex pieces. Symbolically, the requested quantity is
\[
\dim\mathcal P_n(\R^3),\qquad n\in\Z,\quad n>3,
\]
where \(\mathcal P_n(\R^3)\) denotes the intended space of convex partitions.

\textbf{Unresolved formulation qualifications.} The supplied record does not define this space or its dimension. A partition must cover \(\R^3\) by convex pieces, but the record does not specify whether pieces are nonempty, full-dimensional, open, closed, or otherwise restricted. It also does not specify whether pieces are pairwise disjoint as sets, have disjoint interiors with shared boundaries, or use a boundary-assignment convention. No topology, parameterization, or notion of dimension is supplied. Nor is it stated whether pieces are labeled or whether partitions are identified under translations, Euclidean isometries, similarities, or other transformations. Consequently, \(\mathcal P_n(\R^3)\) and \(\dim\) above are placeholders for the source's intended objects; a uniquely defined numerical target cannot be recovered from the supplied record alone.

The supplied background describes the three-dimensional convex-partition realization-space dimension as unresolved beyond a power-diagram lower bound. It supplies neither the numerical lower bound nor the definitions and hypotheses under which that comparison is made. This contextual qualification is retained without asserting a bound or selecting a realization-space convention.
\subsection{Short English statement}
For every integer \(n>3\), determine the dimension of the space of partitions of \(\R^3\) into \(n\) convex pieces. The supplied record leaves the admissible pieces, boundary convention, topology, notion of dimension, and equivalences unspecified, so the intended space requires clarification. Its open status and reference to a power-diagram lower bound are supplied curation claims requiring independent review.
\subsection{Sources}
\begin{itemize}
\item Local UnsolvedMath registry, permanent ID \(30001893\), title: \emph{Dimension of Convex Partition Spaces in Three Dimensions}; source problem code \texttt{catalogue entry}. The supplied record belongs to the Oberwolfach Reports Open Problems collection and records proposed year 2011. No individual proposer is supplied.
\item Supplied source citation: \emph{Discrete Geometry (2012)}, Oberwolfach Reports, pp.\ 2459--2548, DOI \texttt{10.4171/owr/2011/44}. Source URL: \url{https://doi.org/10.4171/owr/2011/44}. The supplied background describes Oberwolfach Reports as documenting mathematical workshops and their open-problem discussions.
\item Supplied literature reference: \emph{OWR 44 (2012)}, \url{https://doi.org/10.4171/owr/2011/44}. Literature checked at: \texttt{2026-08-21}. This date and the assessment that the problem is open beyond a power-diagram lower bound are supplied dated curation, not independently verified here. The background explicitly requests independent verification of time-sensitive status and cited claims.
\item The registry labels the statement \texttt{corrected\_verified}; its supplied verification explanation says only that TeX notation was enclosed in mathematical delimiters without changing the mathematical content. This label does not resolve the missing mathematical conventions.
\end{itemize}
\end{document}
